%-------------------------------------------------------------------------------
% Mattia Giovannini — CV in italiano (versione pubblica), layout Awesome-CV
% Compilare con LuaLaTeX:  latexmk -lualatex CV-IT-public.tex
%-------------------------------------------------------------------------------
\documentclass[11pt, a4paper]{../awesome-cv}
% Shared Awesome-CV settings for the public CVs (IT/EN/FR).
% Each CV loads this after \documentclass{../awesome-cv}; build with LuaLaTeX.

\geometry{left=1.4cm, top=.9cm, right=1.4cm, bottom=1.6cm, footskip=.5cm}

% Same blue accent as the other CVs
\definecolor{awesome}{RGB}{80, 100, 200}
\setbool{acvSectionColorHighlight}{true}
\renewcommand{\acvHeaderSocialSep}{\quad\textbar\quad}

% No hyphenation: words stay whole when application portals extract the text
\hyphenpenalty=10000
\exhyphenpenalty=10000

\name{Mattia}{Giovannini}
\email{mattia.giovannini6@studio.unibo.it}
\homepage{mattia-giovannini.github.io}
\github{frigori}

% Re-load the class fonts with LuaTeX's node renderer: with HarfBuzz, explicit
% hyphens ("eye-tracking") are extracted from the PDF as invisible soft hyphens.
\defaultfontfeatures{Ligatures=TeX, Renderer=Node}
\setmainfont{Source Sans 3}[
  UprightFont=*, ItalicFont=* Italic, BoldFont=* Bold, BoldItalicFont=* Bold Italic,
  FontFace={l}{n}{Font=* Light}, FontFace={l}{it}{Font=* Light Italic},
]
\setsansfont{Source Sans 3}[
  UprightFont=*, ItalicFont=* Italic, BoldFont=* Bold, BoldItalicFont=* Bold Italic,
  FontFace={l}{n}{Font=* Light}, FontFace={l}{it}{Font=* Light Italic},
]
\renewfontfamily\roboto{Roboto}[
  UprightFont=*, ItalicFont=* Italic, BoldFont=* Bold, BoldItalicFont=* Bold Italic,
  FontFace={l}{n}{Font=* Light}, FontFace={l}{it}{Font=* Light Italic},
]


\position{Studente Magistrale in Ingegneria Informatica (Doppio Titolo){\enskip\cdotp\enskip}Laureato in Ingegneria Informatica}
\address{Bologna, Italia}

\begin{document}

\makecvheader[C]
\makecvfooter{}{Mattia Giovannini~~~·~~~Curriculum Vitae}{\thepage}

%-------------------------------------------------------------------------------
\cvsection{Profilo}

\begin{cvparagraph}
Laureato in Ingegneria Informatica presso l'Università di Bologna e iscritto alla Laurea Magistrale in Ingegneria
Informatica presso lo stesso ateneo, in doppio titolo con il Master in Ubiquitous Systems and Networks
dell'Université Côte d'Azur (Polytech Nice Sophia), con un forte interesse per l'intelligenza artificiale, il machine
learning e l'analisi dei dati. Il mio percorso mi ha fornito solide basi in algoritmi, statistica e ingegneria del
software, che ho applicato in un contesto di ricerca durante un semestre Erasmus in Francia. Sono motivato dalla
sfida di progettare sistemi intelligenti per la risoluzione di problemi reali.
\end{cvparagraph}

%-------------------------------------------------------------------------------
\cvsection{Formazione}
\begin{cventries}
  \cventry
    {Doppio titolo con il Master in Ubiquitous Systems and Networks, Université Côte d'Azur}
    {Laurea Magistrale in Ingegneria Informatica — Doppio Titolo}
    {Bologna, Italia \& Nizza, Francia}
    {2026 -- 2028 (In corso)}
    {}

  \cventry
    {Voto di laurea: 94/110}
    {Laurea Triennale in Ingegneria Informatica — Università di Bologna}
    {Bologna, Italia}
    {2021 -- Luglio 2026}
    {
      \begin{cvitems}
        \item {\textbf{Tesi:} \textit{Rilevamento del Disturbo dello Spettro Autistico mediante Machine Learning applicato a Dati di Eye Tracking}, con tecniche di classificazione e analisi dei segnali applicate a registrazioni cliniche di eye-tracking per il supporto allo screening dell'ASD.}
        \item {Corsi rilevanti: Analisi Matematica, Algebra Lineare e Geometria, Probabilità e Statistica, Fondamenti di Informatica, Sistemi Operativi, Reti di Calcolatori, Sistemi Informativi, Ingegneria del Software, Tecnologie Web, Architettura degli Elaboratori.}
      \end{cvitems}
    }

  \cventry
    {Progetto di studio Erasmus}
    {Université de Tours}
    {Tours, Francia}
    {2025}
    {
      \begin{cvitems}
        \item {Sviluppo di un progetto sul \textbf{machine learning per l'analisi del disturbo dello spettro autistico (ASD)}, con classificazione ed estrazione di caratteristiche da dati di eye-tracking.}
        \item {Lavoro in un contesto accademico internazionale e francofono; rafforzamento di analisi dei dati, Python e metodologia scientifica.}
      \end{cvitems}
    }

  \cventry
    {Diploma di Liceo Scientifico}
    {Liceo Augusto Righi}
    {Bologna, Italia}
    {2016 -- 2021}
    {}
\end{cventries}

%-------------------------------------------------------------------------------
\cvsection{Progetti}
\begin{cventries}
  \cventry
    {Progetto GenAI}
    {BOOM GEN AI Challenge 2026 — VEM Sistemi}
    {}
    {2026}
    {
      \begin{cvitems}
        \item {Partecipazione allo sviluppo di un \textbf{assistente basato su GenAI} per l'analisi automatica di bandi di gara, a supporto della valutazione \textbf{GO/NO-GO}.}
        \item {Progettazione e implementazione di un modulo di \textbf{scoring} in \textbf{Python} per valutare la coerenza tra requisiti del bando e capacità aziendali.}
        \item {\textbf{Prompt engineering} e servizi \textbf{Azure OpenAI} per estrazione, classificazione e arricchimento dei requisiti; analisi della pipeline di parsing di bandi non strutturati.}
      \end{cvitems}
    }

  \cventry
    {Project work · Enterprise Platforms and DevOps M, Università di Bologna}
    {Scheduler Kubernetes resource-aware}
    {}
    {2026 (In corso)}
    {
      \begin{cvitems}
        \item {Sviluppo di uno \textbf{scheduler Kubernetes personalizzato} in Python che assegna ogni pod al nodo a latenza minore tra quelli con CPU, memoria e slot sufficienti.}
        \item {Misura della latenza dei nodi in tempo reale, binding tramite API Kubernetes e motivazione di ogni decisione con eventi \textit{Scheduled}/\textit{FailedScheduling}.}
        \item {In corso l'estensione allo \textbf{smart placement e alla sostituzione} di servizi su nodi edge in movimento (robot umanoidi tra postazioni di lavoro).}
      \end{cvitems}
    }

  \cventry
    {Discrete and Network Optimization Problems M, Università di Bologna}
    {Benchmark di ripristino da guasti di link SDN}
    {\href{https://github.com/frigori/net-restoration-bench}{\faGithub\ net-restoration-bench}}
    {2026 (In corso)}
    {
      \begin{cvitems}
        \item {Confronto tra \textbf{backup condiviso (SBPP)}, \textbf{Fast Reroute} e \textbf{ricalcolo globale MILP} su overhead di capacità, tasso di ripristino, tempo di reazione e path stretch.}
        \item {Python con NetworkX e PuLP; un'unica interfaccia plan/react esegue le strategie su topologie e scenari di guasto identici.}
      \end{cvitems}
    }

  \cventry
    {Lean Startup Lab · Technology Entrepreneurship}
    {AgenticOS — Audit multi-agente con IA}
    {\href{https://agenticos-mvp.netlify.app}{\faGlobe\ MVP online}}
    {2026}
    {
      \begin{cvitems}
        \item {Sviluppo di un \textbf{sistema multi-agente Red Team / Blue Team} (Gemini) che analizza perché gli agenti IA di shopping ignorano un prodotto moda.}
        \item {Un agente Blue Team traduce le correzioni in \textbf{JSON-LD} e \textbf{llms.txt} e le valuta in sandbox; rilasciato come MVP con interfaccia web e motore containerizzato.}
      \end{cvitems}
    }

  \cventry
    {Progetto personale}
    {Homelab self-hosted}
    {\href{https://mattia-giovannini.github.io/homelab/}{\faGlobe\ pagina homelab}}
    {In corso}
    {
      \begin{cvitems}
        \item {Gestione di servizi containerizzati su nodi Raspberry Pi con \textbf{Docker, Podman rootless e Kubernetes}, dietro un reverse proxy con TLS.}
        \item {Monitoraggio Prometheus/Grafana, sicurezza con CrowdSec; banco di prova per dare all'IA agentica un \textbf{accesso limitato e verificabile}.}
      \end{cvitems}
    }
\end{cventries}

%-------------------------------------------------------------------------------
\cvsection{Competenze}
\begin{cvskills}
  \cvskill{Programmazione}{C, Java, Python, JavaScript, HTML/CSS, SQL, Bash}
  \cvskill{ML e dati}{scikit-learn, NumPy, pandas, Matplotlib, Jupyter Notebook}
  \cvskill{Sistemi e reti}{Linux, TCP/IP, programmazione socket, database relazionali}
  \cvskill{Strumenti}{Git, \LaTeX, MATLAB, Docker, Podman, Kubernetes}
  \cvskill{Interessi}{Machine learning, analisi dei dati, containerizzazione, automazione di sistema}
\end{cvskills}

%-------------------------------------------------------------------------------
\cvsection{Lingue}
\begin{cvskills}
  \cvskill{Italiano}{Madrelingua}
  \cvskill{Inglese — C1}{Duolingo English Test 130/160 (C1) · CLA C1, Università di Bologna · Cambridge First Certificate B2 (178/190)}
  \cvskill{Francese — C1}{CLA C1, Università di Bologna · Corsi B2/C1 presso l'Université de Tours}
\end{cvskills}

%-------------------------------------------------------------------------------
\cvsection{Altre esperienze}
\begin{cventries}
  \cventry
    {Webmaster ed Editor Studentesco}
    {Liceo Augusto Righi}
    {Bologna, Italia}
    {2020 -- 2021}
    {
      \begin{cvitems}
        \item {Progettazione e gestione delle piattaforme studentesche dell'istituto (\textit{Righi Network} e \textit{Il Sottobanco}), a supporto della presenza online del giornale scolastico.}
      \end{cvitems}
    }

  \cventry
    {Educatore Volontario presso Centro Estivo}
    {Estate Ragazzi}
    {Bologna, Italia}
    {2016 -- 2020}
    {
      \begin{cvitems}
        \item {Gestione di gruppi, lavoro di squadra, comunicazione e responsabilità in attività educative.}
      \end{cvitems}
    }
\end{cventries}

\end{document}
